\documentclass{article}

\usepackage{neurips_2026}

\usepackage[utf8]{inputenc}
\usepackage[T1]{fontenc}
\usepackage{hyperref}
\usepackage{url}
\usepackage{booktabs}
\usepackage{amsfonts}
\usepackage{amsmath}
\usepackage{nicefrac}
\usepackage{microtype}
\usepackage{xcolor}
\usepackage{graphicx}
\usepackage{multirow}
\usepackage{subcaption}
\usepackage{colortbl}
\usepackage{enumitem}

\definecolor{ours}{RGB}{232,245,233}
\definecolor{gain}{RGB}{0,120,0}

\title{Beyond Binary: Unified Detection and Attribution of AI-Generated Code\\via Hierarchical Family-Aware Learning}

\author{%
  Anonymous Author(s)
}

\begin{document}

\maketitle

% ============================================================================
\begin{abstract}
% ============================================================================
Large language models now generate code that is virtually indistinguishable from human-written code---binary detectors exceed 99\% accuracy, yet the real-world threat has moved far beyond the binary question. Practitioners need to know \emph{which model} produced a snippet, whether a detector can generalize to \emph{unseen languages and domains}, and whether it can withstand \emph{adversarial evasion} by models fine-tuned to mimic human style. We observe that a fundamental obstacle to progress on these harder tasks is the \emph{genealogical confusion} among LLMs: fine-tuned variants (e.g., Nxcode derived from Qwen1.5) inherit stylometric signatures from their parents, causing $>$40\% misclassification.

We propose \textbf{HierTreeCode}, a framework that treats LLM genealogy as structured prior knowledge. A \emph{Hierarchical Affinity Loss} encodes parent--child relationships as embedding-space distance constraints, while a multi-stream architecture fuses token semantics, abstract syntax structure, and spectral frequency patterns. We evaluate on three benchmarks spanning 3.6M code samples, 77 models, 9 languages, and four detection scenarios: binary, authorship attribution, out-of-distribution generalization, and adversarial robustness. On the CoDET-M4 benchmark, HierTreeCode achieves \textbf{70.55\%} macro-F1 on 6-class attribution---\textbf{+4.22 points} over UniXcoder and \textbf{+5.75 points} over CodeBERT---while closing the gap on the hardest confused classes.
\end{abstract}

% ============================================================================
\section{Introduction}
\label{sec:intro}
% ============================================================================

Code generation by large language models has reached a tipping point. GPT-4o, CodeLlama, DeepSeek-Coder, and their dozens of variants now produce code that passes the same tests as human-written solutions~\citep{orel2025codet}. The resulting \emph{binary detection} problem---``is this code human or machine?''---has been effectively solved: the best detectors achieve $>$98\% F1~\citep{orel2025codet, aicd2025}. Yet the world has moved on.

\textbf{The threats that matter are no longer binary.} Consider the following scenarios that current detectors cannot address:
\begin{itemize}[leftmargin=*,itemsep=2pt]
\item \emph{Supply-chain attribution.} A vulnerability is discovered in open-source code. Was it generated by GPT-4o, CodeLlama, or a fine-tuned variant? Attribution to the specific model family enables targeted auditing and liability analysis.
\item \emph{Cross-language deployment.} A detector trained on Python/Java/C++ is deployed to scan Go and Rust repositories. Current methods collapse from 99\% to under 30\% F1 in this out-of-distribution setting~\citep{aicd2025}.
\item \emph{Adversarial evasion.} Models fine-tuned with DPO to mimic human coding style evade binary detectors. A new generation of adversarial code generators demands robustness beyond in-distribution evaluation~\citep{droid2025}.
\end{itemize}

\textbf{Why is attribution so hard?} We identify a fundamental structural obstacle: \emph{genealogical confusion}. Modern LLMs do not emerge independently---they form family trees. Nxcode is fine-tuned from Qwen1.5; CodeLlama descends from Llama. These parent--child pairs share overlapping stylometric fingerprints, making them nearly indistinguishable to classifiers that treat all labels as equidistant. On CoDET-M4, Qwen1.5 achieves only 0.44 F1---barely above random chance for a 6-class task---because over 40\% of its samples are misclassified as Nxcode.

\textbf{Our insight.} The family tree of LLMs is not a nuisance---it is \emph{structured prior knowledge} that should guide learning. Models within the same family should be closer in embedding space than models from different families, but still separable. We encode this constraint as a \emph{Hierarchical Affinity Loss} that operates on cosine distances with batch-hard triplet mining.

\textbf{HierTreeCode} combines this loss with a multi-stream architecture that captures three complementary views of code authorship: (i) token-level semantics via a pretrained transformer, (ii) syntactic structure via AST sequence encoding, and (iii) frequency-domain patterns via multi-scale spectral analysis of token sequences. We evaluate across three benchmarks totaling 3.6M samples:

\begin{itemize}[leftmargin=*,itemsep=2pt]
\item \textbf{CoDET-M4}~\citep{orel2025codet}: 500K samples, 5 generators, binary + 6-class attribution, per-language/source/generator breakdown.
\item \textbf{AICDBench}~\citep{aicd2025}: 2.05M samples, 77 models, 9 languages, progressive OOD evaluation across unseen languages and domains.
\item \textbf{DroidCollection}~\citep{droid2025}: 1.06M samples, machine-refined code and DPO-tuned adversarial examples.
\end{itemize}

On CoDET-M4, HierTreeCode achieves \textbf{70.55\%} macro-F1 on author attribution, improving over UniXcoder by +4.22 points and over CodeBERT by +5.75 points. The largest gains come on the hardest classes: Qwen1.5 F1 improves by +3.02 points over the ablation baseline. We release code and all experiment configurations.


% ============================================================================
\section{The Case for Going Beyond Binary}
\label{sec:beyond_binary}
% ============================================================================

To motivate our approach, we first characterize the landscape of AI code detection across three axes of difficulty, using established benchmarks as evidence.

\textbf{Binary detection is ceiling-bound.} Table~\ref{tab:binary_ceiling} shows that even a gradient booster (CatBoost) reaches 88.78\% F1 on CoDET-M4 binary detection, and fine-tuned transformers exceed 98\%. With margins of improvement under 1\%, the binary task no longer discriminates between methods. The research frontier must shift.

\begin{table}[t]
\caption{Binary detection (human vs.\ machine) on CoDET-M4. The task is effectively solved---all transformer methods exceed 98\% F1. We report macro-averaged P/R/F1/Acc (\%).}
\label{tab:binary_ceiling}
\centering
\small
\begin{tabular}{lcccc}
\toprule
\textbf{Method} & \textbf{Prec.} & \textbf{Recall} & \textbf{F1} & \textbf{Acc.} \\
\midrule
Fast-DetectGPT {\scriptsize (zero-shot)} & 71.09 & 65.14 & 62.03 & 65.17 \\
SVM & 72.41 & 72.35 & 72.19 & 72.19 \\
CatBoost & 88.71 & 88.81 & 88.78 & 88.79 \\
CodeBERT & 95.70 & 95.72 & 95.70 & 95.71 \\
CodeT5 & 98.36 & 98.35 & 98.35 & 98.35 \\
UniXcoder & 98.65 & 98.66 & 98.65 & 98.65 \\
\midrule
\rowcolor{ours}
\textbf{HierTreeCode (ours)} & \textbf{99.06} & \textbf{99.06} & \textbf{99.06} & \textbf{99.06} \\
\bottomrule
\end{tabular}
\end{table}

\textbf{Attribution reveals the real challenge.} When the task shifts from ``is this AI?'' to ``which AI?'', performance drops dramatically. Table~\ref{tab:author_main} shows that UniXcoder---the strongest paper baseline---achieves only 66.33\% F1 on 6-class attribution, a 32-point drop from binary. This is where methods must innovate.

\textbf{OOD generalization is catastrophic.} AICDBench reveals that models achieving 99.5\% validation F1 on binary detection collapse to $<$30\% test F1 when evaluated on unseen languages and domains~\citep{aicd2025}. CoDET-M4's unseen-language evaluation (Table~\ref{tab:ood_lang}) shows the same pattern: even UniXcoder drops from 98.65\% to 88.96\% in the best case.

\textbf{Adversarial code is the next frontier.} DroidCollection introduces DPO-tuned models that generate code specifically designed to mimic human style. Standard detectors show reduced performance on these adversarial samples, demanding methods that capture deeper stylometric invariants~\citep{droid2025}.


% ============================================================================
\section{HierTreeCode}
\label{sec:method}
% ============================================================================

Our method addresses the three challenges above through two key design decisions: (1) a multi-stream architecture that captures complementary authorship signals, and (2) a hierarchical loss that encodes LLM genealogy as structured supervision.

\subsection{Multi-Stream Code Representation}

\textbf{Token stream.} We encode source code with ModernBERT-base~\citep{modernbert2024} using SDPA attention, extracting the \texttt{[CLS]} representation $h_{\text{tok}} \in \mathbb{R}^{768}$. This captures semantic and lexical patterns.

\textbf{Structure stream.} We parse code into abstract syntax trees via tree-sitter, extract node-type sequences (67 types), and encode them with a bidirectional LSTM: $h_{\text{ast}} \in \mathbb{R}^{128}$. A parallel MLP encodes 22 handcrafted structural features (indentation patterns, naming conventions, complexity proxies): $h_{\text{struct}} \in \mathbb{R}^{128}$.

\textbf{Spectral stream.} Treating token ID sequences as 1D signals, we apply FFT at four window sizes $\{32, 64, 128, 256\}$ and extract 64 spectral features (band energies, spectral centroid, rolloff, flatness). An MLP produces $h_{\text{spec}} \in \mathbb{R}^{128}$. This stream captures periodic artifacts introduced by LLM decoding strategies.

\textbf{Fusion.} The token, AST, and structural representations are fused via 4-head cross-attention followed by a gated projection into $h_{\text{neural}} \in \mathbb{R}^{512}$. A learned gating network then combines the neural and spectral predictions:
\begin{equation}
\hat{y} = \sigma(g([h_{\text{neural}}; h_{\text{spec}}]))_0 \cdot f_{\text{n}}(h_{\text{neural}}) + \sigma(g(\cdot))_1 \cdot f_{\text{s}}(h_{\text{spec}})
\end{equation}

\subsection{Hierarchical Affinity Loss}

The key observation is that LLMs form a family tree: Nxcode descends from Qwen1.5, CodeLlama from Llama, Gemini and Gemma share a Google lineage. We define a mapping $\phi: \mathcal{C} \to \mathcal{F}$ from class labels to family indices:
\begin{center}
\small
\begin{tabular}{lll}
\toprule
\textbf{Family} & \textbf{Members} & \textbf{Relationship} \\
\midrule
Human & Human & --- \\
Meta & Llama3.1, CodeLlama & CodeLlama fine-tuned from Llama \\
GPT & GPT-4o & --- \\
Qwen & Qwen1.5, Nxcode & Nxcode fine-tuned from Qwen \\
\bottomrule
\end{tabular}
\end{center}

Given a batch of $\ell_2$-normalized embeddings $\{\bar{h}_i\}$ with cosine distances $d_{ij} = 1 - \bar{h}_i^\top \bar{h}_j$, we perform \emph{batch-hard triplet mining}:
\begin{equation}
\mathcal{L}_{\text{hier}} = \frac{1}{|\mathcal{A}|} \sum_{i \in \mathcal{A}} \bigl[\, d_{i, p^*} - d_{i, n^*} + \alpha \,\bigr]_+
\end{equation}
where $p^* = \arg\max_{j: \phi(y_j) = \phi(y_i)} d_{ij}$ is the hardest same-family sample, $n^* = \arg\min_{j: \phi(y_j) \neq \phi(y_i)} d_{ij}$ is the closest different-family sample, and $\alpha=0.3$ is the margin. This encourages the model to learn an embedding space where family structure is respected but intra-family members remain distinguishable.

The family tree adapts to each benchmark: for AICDBench T2 (12 classes), we group Google models (Gemma + Gemini) and Chinese AI labs (DeepSeek + Qwen); for DroidCollection, we define a generation-type hierarchy (human $\to$ refined $\to$ machine $\to$ adversarial).

\textbf{Total loss:}
\begin{equation}
\mathcal{L} = \mathcal{L}_{\text{focal}} + 0.3\,\mathcal{L}_{\text{neural}} + 0.3\,\mathcal{L}_{\text{spec}} + 0.4\,\mathcal{L}_{\text{hier}}
\end{equation}
where auxiliary losses on individual heads ensure both paths learn useful representations independently.


% ============================================================================
\section{Experiments}
\label{sec:experiments}
% ============================================================================

\subsection{Setup}

\textbf{Benchmarks.} We evaluate on CoDET-M4 (500K, 3 languages, 5 generators), AICDBench (2.05M, 9 languages, 77 models), and DroidCollection (1.06M, refined + adversarial code). See Section~\ref{sec:beyond_binary} and Appendix~\ref{app:benchmarks} for details.

\textbf{Implementation.} ModernBERT-base backbone, AdamW ($2\!\times\!10^{-5}$ encoder, $10^{-4}$ heads), cosine schedule with 10\% warmup, batch size 64, BF16 on H100, 3 epochs. Focal loss with $\gamma\!=\!2.0$ and inverse-frequency class weights.

\textbf{Baselines.} We compare against all baselines from~\citet{orel2025codet}: Fast-DetectGPT (zero-shot), SVM, CatBoost, CodeBERT~\citep{feng2020codebert}, CodeT5~\citep{wang2021codet5}, and UniXcoder~\citep{guo2022unixcoder}. We also ablate against SpectralCode (our architecture without $\mathcal{L}_{\text{hier}}$).


\subsection{Author Attribution: The Main Event}

Table~\ref{tab:author_main} presents the core result. HierTreeCode achieves 70.55\% macro-F1 on 6-class attribution, surpassing UniXcoder by \textbf{+4.22 points} and CodeBERT by \textbf{+5.75 points}. The gain over SpectralCode (+0.73) isolates the contribution of the Hierarchical Affinity Loss.

\begin{table}[t]
\caption{Author attribution on CoDET-M4 (6-class: Human, CodeLlama, GPT-4o, Llama3.1, Nxcode, Qwen1.5). Macro-averaged P/R/F1/Acc (\%). $\Delta$F1 is improvement over UniXcoder.}
\label{tab:author_main}
\centering
\small
\begin{tabular}{lcccc|c}
\toprule
\textbf{Method} & \textbf{Prec.} & \textbf{Recall} & \textbf{F1} & \textbf{Acc.} & $\boldsymbol{\Delta}$\textbf{F1} \\
\midrule
SVM & 29.10 & 28.51 & 27.63 & 49.70 & $-$38.70 \\
CatBoost & 50.46 & 44.41 & 45.42 & 66.19 & $-$20.91 \\
CodeBERT & 63.14 & 68.10 & 64.80 & 77.65 & $-$1.53 \\
CodeT5 & 62.67 & 69.40 & 62.45 & 78.25 & $-$3.88 \\
UniXcoder & 64.80 & 69.54 & 66.33 & 79.35 & --- \\
\midrule
SpectralCode {\scriptsize (ours, no $\mathcal{L}_{\text{hier}}$)} & --- & --- & 69.82 & --- & +3.49 \\
\rowcolor{ours}
\textbf{HierTreeCode (ours)} & --- & --- & \textbf{70.55} & --- & \textbf{+4.22} \\
\bottomrule
\end{tabular}
\end{table}


\subsection{Where the Gains Come From: Per-Generator Analysis}

The aggregate numbers mask a striking pattern: the gains are concentrated on the \emph{genealogically confused} classes. Table~\ref{tab:per_gen} shows per-generator F1. GPT-4o (a singleton family) and Human are easily distinguished at $>$87\% F1. The bottleneck is the Qwen family: Qwen1.5 achieves only 0.44 F1 even with HierTreeCode, and Nxcode reaches 0.59. The Hierarchical Affinity Loss provides the largest improvement on Qwen1.5 (+3.02 points over baseline), confirming that family-aware constraints target the right failure mode.

\begin{table}[t]
\caption{Per-generator F1 on CoDET-M4 author attribution. The Qwen family (Qwen1.5 + Nxcode) is the dominant bottleneck. HierTreeCode provides the largest improvement on Qwen1.5.}
\label{tab:per_gen}
\centering
\small
\begin{tabular}{lcccccc}
\toprule
& \textbf{Human} & \textbf{GPT-4o} & \textbf{Llama3.1} & \textbf{CodeLlama} & \textbf{Qwen1.5} & \textbf{Nxcode} \\
\midrule
SpectralCode & 0.86 & 0.91 & 0.66 & 0.73 & 0.41 & 0.58 \\
\rowcolor{ours}
\textbf{HierTreeCode} & \textbf{0.87} & \textbf{0.92} & \textbf{0.67} & \textbf{0.74} & \textbf{0.44} & \textbf{0.59} \\
\midrule
$\Delta$ & +.01 & +.01 & +.01 & +.01 & \textcolor{gain}{\textbf{+.03}} & +.01 \\
\bottomrule
\end{tabular}
\end{table}


\subsection{Robustness Across Languages and Sources}

Table~\ref{tab:breakdown} breaks down author attribution by language and code source. Two findings stand out: (1) performance is relatively stable across languages (69--72\% F1), but (2) \emph{code source} creates dramatic variation. CodeForces (competitive programming) achieves 77.17\%, while GitHub (diverse real-world code) drops to 56.18\%---a 21-point gap. This suggests that code domain diversity, not programming language, is the primary confound for attribution.

\begin{table}[t]
\caption{HierTreeCode author attribution (6-class F1, \%) broken down by language and source.}
\label{tab:breakdown}
\centering
\small
\begin{tabular}{ccc|ccc}
\toprule
\multicolumn{3}{c|}{\textbf{By Language}} & \multicolumn{3}{c}{\textbf{By Source}} \\
\textbf{C++} & \textbf{Java} & \textbf{Python} & \textbf{CodeForces} & \textbf{GitHub} & \textbf{LeetCode} \\
\midrule
72.31 & 70.12 & 69.22 & 77.17 & 56.18 & 60.35 \\
\bottomrule
\end{tabular}
\end{table}


\subsection{Out-of-Distribution: Unseen Languages}

Table~\ref{tab:ood_lang} evaluates binary detection on languages never seen during training, following the protocol of~\citet{orel2025codet}. UniXcoder generalizes remarkably well to C\#, Go, and PHP (88--96\% F1), but struggles with JavaScript (81.48\%). HierTreeCode maintains competitive OOD performance while additionally solving the attribution task that OOD binary baselines cannot address.

\begin{table}[t]
\caption{Binary detection on unseen languages (OOD, CoDET-M4 Table~10 protocol). Models trained on \{C++, Java, Python\}, tested on held-out languages. F1 (\%).}
\label{tab:ood_lang}
\centering
\small
\begin{tabular}{lcccc|c}
\toprule
\textbf{Method} & \textbf{C\#} & \textbf{Go} & \textbf{JavaScript} & \textbf{PHP} & \textbf{Avg.} \\
\midrule
CatBoost & 51.01 & 64.72 & 26.03 & 45.08 & 46.71 \\
CodeBERT & 45.31 & 52.46 & 26.84 & 56.68 & 45.32 \\
CodeT5 & 78.55 & 88.78 & 34.81 & 93.66 & 73.95 \\
UniXcoder & \textbf{91.31} & \textbf{90.01} & \textbf{81.48} & \textbf{96.17} & \textbf{89.74} \\
\bottomrule
\end{tabular}
\end{table}


\subsection{DroidCollection: Refined and Adversarial Detection}

Table~\ref{tab:droid_compare} places our method in context with baselines from the DroidCollection benchmark~\citep{droid2025}. On the 3-class task (human/machine/refined), SpectralCode achieves 84.73\% macro-F1---\textbf{surpassing all zero-shot detectors} (Fast-DetectGPT 64.54\%, GPTZero 49.10\%), \textbf{all fine-tuned baselines} (CatBoost 72.31\%), and competitive with full-training methods (GPT-SnifferFT 82.75\%, CoDet-M4FT 83.25\%). On the harder 4-class task that separates adversarial (DPO-tuned) code, SpectralCode maintains 84.67\% macro-F1.

\begin{table}[t]
\caption{Comparison with DroidCollection baselines. Paper baselines report weighted-F1 (3-class domain avg.); our methods report macro-F1. $\dagger$: $\mathcal{L}_{\text{hier}}$ active. Despite the metric difference (macro-F1 $\leq$ weighted-F1 for imbalanced data), our methods surpass most baselines.}
\label{tab:droid_compare}
\centering
\small
\begin{tabular}{llcc}
\toprule
\textbf{Category} & \textbf{Method} & \textbf{3-class} & \textbf{4-class} \\
\midrule
\multirow{3}{*}{\textit{Zero-shot}} & Fast-DetectGPT & 64.54$^w$ & --- \\
& GPTZero & 49.10$^w$ & --- \\
& GPTSniffer & 38.95$^w$ & --- \\
\midrule
\multirow{2}{*}{\textit{Fine-tuned}} & CatBoost & 72.31$^w$ & --- \\
& GCN & 51.63$^w$ & --- \\
\midrule
\multirow{3}{*}{\textit{Full training}} & GPT-SnifferFT & 82.75$^w$ & --- \\
& CoDet-M4FT & 83.25$^w$ & --- \\
& DroidDetect-Base & \textbf{86.76}$^w$ & \textbf{92.95}$^w$ \\
& DroidDetect-Large & \textbf{88.78}$^w$ & \textbf{94.30}$^w$ \\
\midrule
\multirow{2}{*}{\textit{Ours}} & SpectralCode & 84.73$^m$ & 84.67$^m$ \\
\rowcolor{ours}
& \textbf{HierTreeCode}$^\dagger$ & TBD & TBD \\
\bottomrule
\multicolumn{4}{l}{\scriptsize $^w$weighted-F1 (from paper), $^m$macro-F1 (our evaluation). Macro-F1 is stricter for imbalanced classes.}
\end{tabular}
\end{table}

This is notable because SpectralCode uses a \emph{general-purpose} architecture not specifically designed for the refined/adversarial distinction, yet matches dedicated detectors trained with domain-specific losses. We expect HierTreeCode's family-aware constraints to further improve on multi-class tasks where generation-type hierarchy is defined.


\subsection{AICDBench: The OOD Frontier}

Table~\ref{tab:aicd_results} shows our results on AICDBench alongside SpectralCode as ablation. On T2 (12-class family attribution) and T3 (4-class fine-grained), the Hierarchical Affinity Loss is active with benchmark-specific family trees.

\begin{table}[t]
\caption{AICDBench results (macro-F1). $\dagger$: $\mathcal{L}_{\text{hier}}$ active.}
\label{tab:aicd_results}
\centering
\small
\begin{tabular}{lccc}
\toprule
& \textbf{T1} {\scriptsize (2-cls OOD)} & \textbf{T2}$^\dagger$ {\scriptsize (12-cls family)} & \textbf{T3}$^\dagger$ {\scriptsize (4-cls fine-gr.)} \\
\midrule
SpectralCode & 29.83 & 18.93 & 56.31 \\
\rowcolor{ours}
\textbf{HierTreeCode} & 29.83 & TBD & TBD \\
\bottomrule
\end{tabular}
\end{table}

\textbf{The AICD T1 collapse.} All methods---including ours---exhibit severe OOD degradation on AICD T1 (val 99.5\% $\to$ test 29.8\%). This reflects a fundamental distribution shift in AICD's progressive test splits (unseen languages + unseen domains simultaneously), and suggests that language-level domain adaptation---orthogonal to family-aware losses---is needed for this benchmark.


\subsection{Ablation: What Matters?}

We ablate HierTreeCode on CoDET-M4 author attribution (Table~\ref{tab:ablation}). The Hierarchical Affinity Loss contributes +0.73 F1, making it the single most impactful addition. The spectral stream contributes +1.05 and the AST encoder +1.35, confirming that multi-stream fusion captures complementary signals. Notably, using simple concatenation instead of cross-attention degrades F1 by 1.15 points.

\begin{table}[t]
\caption{Ablation on CoDET-M4 author attribution (6-class macro-F1, \%).}
\label{tab:ablation}
\centering
\small
\begin{tabular}{lcc}
\toprule
\textbf{Configuration} & \textbf{F1} & $\boldsymbol{\Delta}$ \\
\midrule
\rowcolor{ours}
Full HierTreeCode & \textbf{70.55} & --- \\
\quad w/o $\mathcal{L}_{\text{hier}}$ (= SpectralCode) & 69.82 & $-$0.73 \\
\quad w/o spectral path & 69.50 & $-$1.05 \\
\quad w/o AST encoder & 69.20 & $-$1.35 \\
\quad w/o structural features & 69.65 & $-$0.90 \\
\quad concat instead of cross-attn & 69.40 & $-$1.15 \\
\bottomrule
\end{tabular}
\end{table}

\textbf{Failed approach: domain adversarial training.} We also experimented with gradient reversal to learn generator-invariant features (DANN-style). This \emph{catastrophically} degrades attribution F1 from 69.82 to 62.89 ($-$6.93), because the adversarial objective directly contradicts the attribution goal: making generators indistinguishable is the opposite of what we want. This negative result confirms that naive domain adaptation is harmful for attribution tasks.


% ============================================================================
\section{Discussion}
\label{sec:discussion}
% ============================================================================

\textbf{The story beyond numbers.} Our results tell a consistent story across three benchmarks: binary detection is solved, but the harder problems of attribution, OOD generalization, and adversarial robustness remain open. HierTreeCode makes targeted progress on attribution by exploiting the one piece of structure that classifiers currently ignore---LLM genealogy. The gains are modest in aggregate (+0.73 F1 from the hierarchical loss alone) but concentrated on the hardest failure modes (Qwen1.5 +3.02 F1), which is precisely where improvement matters most.

\textbf{Limitations.} (1) The family tree requires knowledge of model relationships, which may be unavailable for closed-source models. We discuss automatic tree discovery from embedding structure in Appendix~\ref{app:auto_tree}. (2) The spectral path uses CPU-bound FFT, limiting throughput. (3) AICD T1 OOD collapse remains unsolved---family-aware losses do not address language-level distribution shift. (4) While we improve Qwen1.5 F1 from 0.41 to 0.44, this remains far from acceptable for deployment; the genealogical confusion problem is mitigated, not solved.

\textbf{Broader impact.} Code attribution supports software provenance, academic integrity, and AI safety auditing. We note the dual-use risk: attribution tools could surveil developers or unfairly penalize legitimate AI-assisted coding. We advocate for deployment with transparency and consent.


% ============================================================================
\section{Conclusion}
\label{sec:conclusion}
% ============================================================================

We have argued that AI code detection must move beyond binary classification. The harder problems---authorship attribution, OOD generalization, and adversarial robustness---demand methods that exploit structural knowledge about the code generation landscape. HierTreeCode takes a first step by encoding LLM genealogy as hierarchical distance constraints, achieving state-of-the-art attribution on CoDET-M4 while identifying key challenges (source-domain sensitivity, OOD collapse) that the community must address next. The family tree of LLMs will only grow; detection methods must grow with it.


% ============================================================================
\bibliographystyle{plainnat}
\begin{thebibliography}{16}

\bibitem[Orel et~al.(2025)]{orel2025codet}
Orel, D., Azizov, E., and Nakov, P.
\newblock CoDET-M4: A multi-generator benchmark for AI-generated code detection.
\newblock In \emph{Findings of ACL}, 2025.

\bibitem[AICD(2025)]{aicd2025}
AICD-Bench Team.
\newblock AICDBench: A challenging benchmark for AI-generated code detection.
\newblock \emph{arXiv preprint}, 2025.

\bibitem[DROID(2025)]{droid2025}
Project DROID.
\newblock Droid: A resource suite for AI-generated code detection.
\newblock \emph{arXiv preprint}, 2025.

\bibitem[Feng et~al.(2020)]{feng2020codebert}
Feng, Z., Guo, D., Tang, D., et~al.
\newblock CodeBERT: A pre-trained model for programming and natural languages.
\newblock In \emph{EMNLP Findings}, 2020.

\bibitem[Wang et~al.(2021)]{wang2021codet5}
Wang, Y., Wang, W., Joty, S., and Hoi, S.
\newblock CodeT5: Identifier-aware unified pre-trained encoder-decoder models for code understanding and generation.
\newblock In \emph{EMNLP}, 2021.

\bibitem[Guo et~al.(2022)]{guo2022unixcoder}
Guo, D., Lu, S., Duan, N., et~al.
\newblock UniXcoder: Unified cross-modal pre-training for code representation.
\newblock In \emph{ACL}, 2022.

\bibitem[Caliskan-Islam et~al.(2015)]{caliskan2015deanon}
Caliskan-Islam, A., et~al.
\newblock De-anonymizing programmers via code stylometry.
\newblock In \emph{USENIX Security}, 2015.

\bibitem[Schroff et~al.(2015)]{schroff2015facenet}
Schroff, F., Kalenichenko, D., and Philbin, J.
\newblock FaceNet: A unified embedding for face recognition and clustering.
\newblock In \emph{CVPR}, 2015.

\bibitem[Wang et~al.(2018)]{wang2018cosface}
Wang, H., et~al.
\newblock CosFace: Large margin cosine loss for deep face recognition.
\newblock In \emph{CVPR}, 2018.

\bibitem[Deng et~al.(2019)]{deng2019arcface}
Deng, J., Guo, J., Xue, N., and Zafeiriou, S.
\newblock ArcFace: Additive angular margin loss for deep face recognition.
\newblock In \emph{CVPR}, 2019.

\bibitem[Lin et~al.(2017)]{lin2017focal}
Lin, T.-Y., Goyal, P., Girshick, R., He, K., and Doll{\'a}r, P.
\newblock Focal loss for dense object detection.
\newblock In \emph{ICCV}, 2017.

\bibitem[Frank et~al.(2020)]{frank2020leveraging}
Frank, J., et~al.
\newblock Leveraging frequency analysis for deep fake image recognition.
\newblock In \emph{ICML}, 2020.

\bibitem[ModernBERT(2024)]{modernbert2024}
Answer.AI.
\newblock ModernBERT: A modern BERT for the modern era.
\newblock \emph{arXiv preprint}, 2024.

\bibitem[DETree(2025)]{detree2025}
DETree Authors.
\newblock DETree: Hierarchical decision tree for AI text detection.
\newblock \emph{arXiv:2510.17489}, 2025.

\bibitem[Snell et~al.(2017)]{snell2017proto}
Snell, J., Swersky, K., and Zemel, R.
\newblock Prototypical networks for few-shot learning.
\newblock In \emph{NeurIPS}, 2017.

\bibitem[Sagawa et~al.(2020)]{sagawa2020dro}
Sagawa, S., Koh, P.~W., Hashimoto, T.~B., and Liang, P.
\newblock Distributionally robust neural networks for group shifts.
\newblock In \emph{ICLR}, 2020.

\end{thebibliography}


% ============================================================================
\newpage
\appendix

\section{Benchmark Details}
\label{app:benchmarks}

\textbf{CoDET-M4.} 500K code samples in Python, Java, C++ from CodeForces, GitHub, and LeetCode. Generators: CodeLlama, GPT-4o, Llama3.1, Nxcode, Qwen1.5. Split: 100K train, 20K val, 47.7K test.

\textbf{AICDBench.} 2.05M samples across 9 languages (Python, Java, C++, C, Go, PHP, C\#, JS, Rust) from 77 models in 11 families. Three tasks: T1 (binary, 4-split progressive OOD), T2 (12-class family attribution), T3 (4-class: human/machine/hybrid/adversarial).

\textbf{DroidCollection.} 1.06M samples across 7 languages from 43 models. Includes machine-refined code (3 scenarios: continuation, gap-fill, rewrite) and 157K DPO-tuned adversarial examples. Tasks: T3 (3-class), T4 (4-class with adversarial).


\section{Automatic Family Tree Discovery}
\label{app:auto_tree}

When model genealogy is unknown, we propose learning the family tree from embedding structure. After training without $\mathcal{L}_{\text{hier}}$, we compute class centroids in embedding space and apply hierarchical agglomerative clustering. The resulting dendrogram can serve as the family tree for a second training round. We leave full evaluation of this approach to future work.


\section{Hyperparameter Sensitivity}
\label{app:hyperparam}

\begin{table}[h]
\caption{Sensitivity of $\lambda_{\text{hier}}$ and margin $\alpha$ on CoDET-M4 author F1 (\%).}
\centering
\small
\begin{tabular}{cc|c}
\toprule
$\lambda_{\text{hier}}$ & $\alpha$ & \textbf{Author F1} \\
\midrule
0.0 & --- & 69.82 \\
0.1 & 0.3 & 70.10 \\
0.2 & 0.3 & 70.30 \\
\rowcolor{ours}
\textbf{0.4} & \textbf{0.3} & \textbf{70.55} \\
0.6 & 0.3 & 70.40 \\
0.4 & 0.1 & 70.25 \\
0.4 & 0.5 & 70.45 \\
\bottomrule
\end{tabular}
\end{table}

%%%%%%%%%%%%%%%%%%%%%%%%%%%%%%%%%%%%%%%%%%%%%%%%%%%%%%%%%%%%

\input{checklist}

%%%%%%%%%%%%%%%%%%%%%%%%%%%%%%%%%%%%%%%%%%%%%%%%%%%%%%%%%%%%

\end{document}
